%% Szkielet wygenerowany z docs/02_spis_tresci.md (plan v1, 11.11.2026).
%% Blok planroboczy usuń, gdy podrozdział jest napisany.

\chapter*{Dodatek A. Pełne treści promptów}
\addcontentsline{toc}{chapter}{Dodatek A. Pełne treści promptów}
\label{app:a}

\begin{planroboczy}
Prompty wstawiane z~plików \texttt{code/prompts/} (np. \texttt{\textbackslash lstinputlisting}), żeby praca pokazywała dokładnie to, co dostał model.
\end{planroboczy}

\cleardoublepage
\chapter*{Dodatek B. Parametry środowiska i~kategorie stanów}
\addcontentsline{toc}{chapter}{Dodatek B. Parametry środowiska i~kategorie stanów}
\label{app:b}

\begin{planroboczy}
Tabele parametrów środowiska ogrodowego i~granic kategorii stanów.
\end{planroboczy}

\cleardoublepage
\chapter*{Dodatek C. Instrukcja odtworzenia wyników}
\addcontentsline{toc}{chapter}{Dodatek C. Instrukcja odtworzenia wyników}
\label{app:c}

\begin{planroboczy}
Repozytorium, wersje oprogramowania i~modelu, polecenia.
\end{planroboczy}

\cleardoublepage
\chapter*{Dodatek D. Wyniki uzupełniające}
\addcontentsline{toc}{chapter}{Dodatek D. Wyniki uzupełniające}
\label{app:d}

\begin{planroboczy}
Krzywe psychometryczne dla wszystkich wariantów promptu.
\end{planroboczy}
